% Source: https://artofproblemsolving.com/community/c3h2324875_geometry_problem
% Stable identity: aops-topic-2324875-q01
Use position vectors. Since $R$ is on $CA$ with $RA=\dfrac13CA$, and $Q$ is on $BC$ with $BQ=\dfrac23BC$,
\[R=\dfrac{2A+C}{3},\qquad Q=\dfrac{B+2C}{3}.\]
Let $X$ be the reflection of $Q$ in $R$, that is, $X=2R-Q$, so that $R$ is the midpoint of $QX$. Then
\[X=\dfrac{4A+2C-B-2C}{3}=\dfrac{4A-B}{3}=A+\dfrac{A-B}{3}.\]
So $X$ lies on line $AB$, beyond $A$ (on the opposite side of $A$ from $B$), with $AX=\dfrac13AB$.

\begin{asy}
usepackage("amsmath");
size(7.5cm);
defaultpen(fontsize(10pt));
pair A=(0,0), B=(6,0), C=(1,4);
pair R=(2*A+C)/3, Q=(B+2*C)/3, X=2*R-Q;
pair P=(23/21,0);
draw(A--B--C--cycle,linewidth(0.9));
draw(X--A,dashed);
draw(P--R,blue); draw(P--Q,blue); draw(R--Q,red); draw(R--X,red+dashed); draw(P--X,blue+dashed);
pair u=unit(Q-R), v=unit(P-R);
draw(R+0.25*u--R+0.25*u+0.25*v--R+0.25*v);
dot(A); dot(B); dot(C); dot(P); dot(Q); dot(R); dot(X);
label("$A$",A,S); label("$B$",B,SE); label("$C$",C,N); label("$P$",P,S);
label("$Q$",Q,NE); label("$R$",R,NW); label("$X$",X,S);
\end{asy}

Since $\angle PRQ=90^\circ$, the line $PR$ is perpendicular to $QX$ at its midpoint $R$: it is the perpendicular bisector of $QX$. (Here $P\ne R$, because $P$ is on $AB$, $R$ is on $CA$, and $R\ne A$.) Hence $PQ=PX$, and in the isosceles triangle $PQX$ the segment $PR$ from the apex to the midpoint of the base bisects the apex angle:
\[\angle RPQ=\angle RPX.\]
Finally, $P$ lies on the segment $AB$ (and $P\ne A$), while $X$ lies on the extension of $BA$ beyond $A$; so seen from $P$, the points $A$ and $X$ are in the same direction, and the rays $PA$ and $PX$ coincide. Therefore
\[\angle APR=\angle XPR=\angle RPQ.\qquad\blacksquare\]
